\documentclass[11pt,a4paper]{article}
\usepackage[margin=25mm]{geometry}
\usepackage{amsmath,amssymb,amsthm,booktabs,hyperref}
\newtheorem{theorem}{Theorem}
\newtheorem{lemma}[theorem]{Lemma}
\newtheorem{proposition}[theorem]{Proposition}
\newcommand{\Sz}{\operatorname{Sz}}
\newcommand{\Sp}{\operatorname{Sp}}
\newcommand{\F}{\mathbb F}
\newcommand{\gen}[1]{\langle#1\rangle}
\title{Symmetric matrices in the Suzuki group\\and the scope of a rank argument}
\author{Human Compression project}
\date{}
\begin{document}\maketitle
\begin{abstract}
The proposed classification of symmetric Suzuki matrices in
Spiga's four-page preprint imposes a product identity outside its
stated domain. We identify the missing defect vector and give the
correct matrix criterion. We also verify all six algebraic cases of
the source: they classify the stronger system that was actually
solved, rather than all symmetric Suzuki matrices. The final
non-generation argument is valid with this additional hypothesis.
Together with the previously accepted rank-four counterexample,
reproduced unchanged in the accompanying volume, these results
dispose of the mathematical claims in the specified source edition.
\end{abstract}

\section{The source and the retained results}
We examine Spiga's preprint \emph{On the rank of Suzuki polytopes:
an answer to Hubard and Leemans} \cite{Spiga}, in its edition
arXiv:2010.04080v1 of 8 October 2020.
This is the four-page source registered as F003.
Theorem~1.1 asserts that a chiral polytope with almost simple Suzuki
automorphism group has rank three. The accepted construction
\cite{Counterexample} gives a chiral polytope of type
$\{31,31,25\}$ with automorphism group $\Sz(32768)$ and therefore
refutes that theorem. Its five-page proof, exact field certificate,
and mathematical review are imported without alteration. No search
or acquisition underlying that construction is repeated here.

The source has only one numbered theorem. Its remaining mathematical
content consists of the reduction at the start of Section~3, a matrix
calculation, six cases solving the resulting equations, and the
concluding generation argument. The present note treats each of these.
It does not assert that a later printing has the same wording.

We retain Wilson's characterization in precisely the form quoted
on page~2 of the source: a matrix in $\Sp_4(q)$ belongs to $\Sz(q)$
if and only if it preserves the displayed semilinear product on
\emph{orthogonal pairs}. We also retain Suzuki's standard group
structure: simplicity for $q\geq8$, the single class of involutions,
and the ovoid description in which each involution fixes exactly one
of $q^2+1$ points, with $q-1$ involutions fixing each point. These are
attributed inputs from \cite{Wilson,Suzuki}, not conclusions of our
finite checks. The reduction to the simple group and rank four in
the first paragraph of source Section~3 is expressly attributed
there to Hubard and Leemans \cite[Theorem~1]{HL}. We retain that
quoted implication as a premise; the full chapter is outside the
accessible corpus (historical identifier L004), and is not credited
as recovered here. Our counterexample already lies in the reduced
case, so its refutation does not depend on this reduction.

\section{The symmetric set and the defect vector}
Put $q=2^{2e+1}$ and $t=2^e$, where $e\geq1$, so that $2t^2=q$.
This agrees with the calculation on source page~4. On page~1 the
additional equality $t=\sqrt{2q}$ and the description of $\F_q^4$
as a vector space over $\F_4$ are typographical errors; the consistent
values are $t=\sqrt{q/2}$ and base field $\F_q$.
Let $J$ be the antidiagonal matrix with all four antidiagonal entries
equal to one, and put $f(u,v)=u^{\mathsf T}Jv$ on $V=\F_q^4$.
The product in the source is alternating, additive in both variables,
and satisfies
\[
 (\lambda u)\mathbin\bullet(\mu v)
   =\lambda^t\mu^t(u\mathbin\bullet v).
\]
Its nonzero values on unordered distinct basis pairs are
\[
 e_1\bullet e_2=e_2,\quad e_1\bullet e_3=e_4,
 \quad e_2\bullet e_4=e_1,\quad e_3\bullet e_4=e_3.
 \tag{1}
\]
Both remaining distinct basis products are zero. In particular,
$f(e_1,e_4)=f(e_2,e_3)=1$, although both products are zero.

\begin{lemma}\label{translate}
Let $K$ be a group and let $j\in K$ satisfy $j^2=1$. The map $x\mapsto xj$ is a
bijection
\[
 \{x\in K:xjx=j\}\longrightarrow\{h\in K:h^2=1\}.
\]
If $K\leq\Sp_4(q)$ and $j=J$, its domain consists exactly of the
symmetric matrices in $K$.
\end{lemma}
\begin{proof}
The equation $xjx=j$ is equivalent to $(xj)^2=1$, and the inverse
map is $h\mapsto hj$. For a symplectic matrix,
$x^{\mathsf T}Jx=J$. Comparing this with $xJx=J$, and cancelling
the invertible matrix $Jx$, gives $x^{\mathsf T}=x$. Conversely a
symmetric symplectic matrix satisfies $xJx=J$.
\end{proof}

Thus, writing $X=\{x\in\Sz(q):xJx=J\}$, the retained involution
counts give
\[
 |X|=1+(q^2+1)(q-1)=q^3-q^2+q.
 \tag{2}
\]
The source's proposed list has only $q$ elements. This elementary
discrepancy already occurs in the accepted counterexample paper;
it diagnoses the matrix classification independently of the polytope
intersection and chirality conditions.

\begin{lemma}\label{defect}
Suppose that $D:V\times V\to V$ is alternating and additive in
each variable, with
$D(\lambda u,\mu v)=\lambda^t\mu^tD(u,v)$.
It vanishes on all orthogonal pairs if and only if there is a vector
$r\in V$ such that
\[
 D(u,v)=f(u,v)^t r\qquad(u,v\in V).
 \tag{3}
\]
\end{lemma}
\begin{proof}
The orthogonal basis pairs force $D_{12}=D_{13}=D_{24}=D_{34}=0$,
where $D_{ij}=D(e_i,e_j)$. The pair
$(e_1+e_2,e_3+e_4)$ is orthogonal and gives $D_{14}=D_{23}$.
Denote this common value by $r$. Expanding in the basis yields
\[
 D(u,v)=(u_1v_4+u_4v_1+u_2v_3+u_3v_2)^t r,
\]
which is (3). Conversely (3) vanishes whenever $f(u,v)=0$.
\end{proof}

For $g\in\Sp_4(q)$ take
$D_g(u,v)=gu\bullet gv-g(u\bullet v)$.
It has the hypotheses of Lemma~\ref{defect}. Consequently Wilson's
condition says that the four orthogonal basis defects vanish and
the two other defects are equal. It does not say that the latter
defects vanish separately.

\section{The equations that are valid}
For ease of comparison set
\[
 x=\begin{pmatrix}
 A&B&C&D\\B&E&F&G\\C&F&H&I\\D&G&I&K
 \end{pmatrix}.
 \tag{4}
\]
The letter $K$ here denotes an entry, not a group. The four
orthogonal basis pairs in (1) give exactly the following ten
distinct scalar equations. The numbers are the source numbers.
\begin{align}
 (BG+DE)^t&=B,\tag{3s}\\
 (AE+B^2)^t&=E,\tag{4s}\\
 (CG+DF)^t&=F,\tag{5s}\\
 (AF+BC)^t&=G,\tag{6s}\\
 (BI+DF)^t&=D,\tag{7s}\\
 (CI+DH)^t&=I,\tag{8s}\\
 (AH+C^2)^t&=K,\tag{9s}\\
 (EK+G^2)^t&=A,\tag{10s}\\
 (FK+GI)^t&=C,\tag{11s}\\
 (HK+I^2)^t&=H.\tag{12s}
\end{align}
For example, the four coordinates of $xe_1\bullet xe_2=xe_2$
give (3s)--(6s). The pair $(e_1,e_3)$ gives (7s), (6s), (8s),
(9s); the pair $(e_2,e_4)$ gives (10s), (3s), (11s), (7s);
and $(e_3,e_4)$ gives (11s), (5s), (12s), (8s).
This also specifies the omitted computations on source page~3.

The two nonorthogonal defects, after applying the inverse of the
Frobenius map $z\mapsto z^t$ coordinatewise, are
\[
 \begin{pmatrix}BK+DG\\AG+BD\\CK+DI\\AI+CD\end{pmatrix},
 \qquad
 \begin{pmatrix}EI+FG\\BF+CE\\FI+GH\\BH+CF\end{pmatrix}.
 \tag{5}
\]
Their equality consists of four off-diagonal equations of $xJx=J$.
For instance the second equality is the $(1,2)$ entry, and the
other three are the $(2,4)$, $(3,4)$ and $(1,3)$ entries.
Thus a symmetric symplectic matrix already has equal defects.

\begin{proposition}\label{criterion}
A symmetric matrix $x$ belongs to $\Sz(q)$ if and only if
$xJx=J$ and (3s)--(12s) hold.
\end{proposition}
\begin{proof}
Necessity follows from the form and the four orthogonal basis
products. Conversely $xJx=J$ implies invertibility and, by symmetry,
membership in $\Sp_4(q)$. Equations (3s)--(12s) make the four
orthogonal basis defects zero. Equation (5) and $xJx=J$ make the
remaining two equal. Lemma~\ref{defect} and Wilson's retained
characterization now give $x\in\Sz(q)$.
\end{proof}

Source equations (13)--(20) instead require both vectors in (5)
to be zero. That is a stronger condition and is the substantive
error. The repeated $xe_1\bullet xe_3$ immediately before
$x(e_2\bullet e_4)$ on page~3 is a separate index error: the
intended pair is $(e_2,e_4)$. Equations (21)--(22) are valid:
\[
 D^2+CG+BI+AK=1,\qquad CG+F^2+EH+BI=1.
 \tag{6}
\]
They are the $(1,4)$ and $(2,3)$ entries of $xJx=J$.

\begin{proposition}\label{binary}
For every allowed $q$, the matrix
\[
 x_0=\begin{pmatrix}
 0&1&0&1\\1&1&1&0\\0&1&0&0\\1&0&0&0
 \end{pmatrix}
 \tag{7}
\]
belongs to $X$ but violates source equation (14).
\end{proposition}
\begin{proof}
Direct multiplication gives $x_0^{\mathsf T}=x_0$ and
$x_0Jx_0=J$. The four orthogonal basis defects are zero; both
other defects equal $e_2$. These are calculations over $\F_2$,
so the exponent $t$ changes none of their entries. Lemma~\ref{defect}
proves preservation on every orthogonal pair, hence membership in
the source Suzuki model. But source equation (14), $BD=AG$, reads
$1=0$ for this matrix. This refutes the intermediate classification
uniformly; it is not itself a polytope marking.
\end{proof}

\section{What the six source cases prove}
Let $Z$ be the set of symmetric matrices (4) satisfying (3s)--(12s),
(6), and the eight additional equations
\begin{alignat*}{4}
 DG&=BK,&\quad BD&=AG,&\quad CK&=DI,&\quad AI&=CD,\\
 EI&=GF,& BF&=EC,& FI&=GH,& BH&=CF.
\end{alignat*}
These are source equations (13)--(20), in order. Define
\[
 T=\{\operatorname{diag}(a,a^{2t+1},a^{-2t-1},a^{-1}):
 a\in\F_q^*\}.
 \tag{8}
\]

\begin{theorem}\label{strong}
The solutions of the stronger system are exactly $Z=\{J\}\cup T$.
Thus all six source cases are valid as a classification of $Z$.
This set is a proper subset of $X$.
\end{theorem}
\begin{proof}
We retain the order of the cases and record the equations used.

Suppose first that $G\ne0$. Equations (13), (14), (17), (19) give
\[
 D=BK/G,\quad A=B^2K/G^2,\quad
 F=EI/G,\quad H=EI^2/G^2.
 \tag{9}
\]
In the first equation of (6), the terms $D^2$ and $AK$ cancel,
leaving $CG+BI=1$. Equation (15) gives $CK=BIK/G$.
Multiplication of $CG+BI=1$ by $K/G$ therefore gives $K/G=0$,
so $K=D=A=0$. Similarly (18) and (9) give $EC=BEI/G$;
multiplication by $E/G$ gives $E/G=0$. Equation (4s) now gives
$B=0$, and (9) gives $H=0$. Equation (9s), with $K=0$, gives
$C=0$. This contradicts $CG+BI=1$.

Suppose $G=0$ and $B\ne0$. Equations (13)--(14) give $K=D=0$.
Equation (3s) then gives $B=0$, a contradiction. This is source
Case~3.2; its intervening consequence $BI=1$ from (6) is correct
but is not needed for the contradiction.

Suppose $G=B=0$ and $I\ne0$. Equation (17) gives $E=0$.
The second equation of (6) gives $F^2=1$, whence $F=1$ in
characteristic two. Equation (19) gives $FI=0$, a contradiction.

Suppose $G=B=I=0$ and $C\ne0$. Equation (15) gives $K=0$.
The first equation of (6) gives $D=1$. Equation (16) then gives
$CD=0$, a contradiction.

Suppose $G=B=I=C=0$ and $D\ne0$. Equation (3s) gives $E=0$;
the second equation of (6) gives $F=1$. Equation (6s) gives
$A=0$, and the first equation of (6) gives $D=1$. Equation (8s)
gives $H=0$, and (9s) gives $K=0$. The resulting matrix is $J$.

Finally suppose $G=B=I=C=D=0$. The first equation of (6) gives
$AK=1$, so $A\ne0$ and $K=A^{-1}$. Equation (5s) gives $F=0$;
the second equation of (6) gives $EH=1$. Equation (10s) becomes
$E^t=A^{t+1}$. Raising to the power $2t$, using $2t^2=q$ and
$z^q=z$ in $\F_q$, gives $E=A^{2t+1}$. Thus $H=A^{-2t-1}$,
and the matrix belongs to $T$. These six cases exhaust all choices
of the five entries tested successively.

For the converse, $J$ satisfies every displayed equation by
substitution. A diagonal member of $T$ has $AK=EH=1$.
The only nonzero product equations are (4s), (9s), (10s), (12s).
For example
\[
 (AE)^t=a^{2t^2+2t}=a^{q+2t}=a^{2t+1}=E,
 \qquad (AH)^t=a^{-2t^2}=a^{-1}=K.
\]
The other two follow by inversion. Hence each listed matrix lies
in $Z$. Each is also symplectic, so Proposition~\ref{criterion}
shows $Z\subseteq X$. Proposition~\ref{binary} shows that the
inclusion is proper.
\end{proof}

\section{The generation argument with its exact hypothesis}
We first retain the valid normalization preceding the source
calculation. Suppose $\sigma_1\sigma_2$, $\sigma_2\sigma_3$,
and $\sigma_1\sigma_2\sigma_3$ are involutions in $\Sz(q)$.
Conjugacy of involutions allows the last product to be $J$.
Then $\sigma_1\sigma_2=J\sigma_3^{-1}$ and
$\sigma_2\sigma_3=\sigma_1^{-1}J$. Squaring these products gives
\[
 \sigma_1^{-1},\sigma_3^{-1}\in X.
 \tag{10}
\]
Every step through (10) is valid. Replacing $X$ by $Z$ is not.

\begin{proposition}\label{nongeneration}
Under these rotation relations and this normalization, if
$\sigma_1^{-1},\sigma_3^{-1}\in Z$, the three rotations do not
generate $\Sz(q)$.
\end{proposition}
\begin{proof}
If $\sigma_1^{-1}=J$, then $\sigma_2\sigma_3=1$; if
$\sigma_3^{-1}=J$, then $\sigma_1\sigma_2=1$. Both contradict
the requirement that the adjacent products have order exactly two.
Theorem~\ref{strong} therefore places both endpoint rotations in
$T$. This is an abelian subgroup of order $q-1$, and direct
conjugation by $J$ inverts every element. Consequently
\[
 \gen{\sigma_1,\sigma_2,\sigma_3}
 =\gen{\sigma_1,\sigma_3,J}\leq T\rtimes\gen J.
\]
The containing group has order $2(q-1)$ and is soluble: its
commutator subgroup is contained in $T$. It cannot equal the
nonabelian simple group $\Sz(q)$.
\end{proof}

This recovers source Section~3.7 as a conditional statement.
The accepted rank-four marking satisfies the genuine relations
and the polytope intersection property; after normalization its
endpoint matrix fails equation (13) of the stronger system.
The missing hypothesis is therefore a genuine restriction on the
construction, not an alternative coordinate convention.

\section{Attribution and disposition}
The equations from orthogonal pairs, the six successive algebraic
cases, and the torus non-generation argument originate in the source.
We have retained them with the correct hypotheses and supplied the
omitted coordinate checks and the converse to the six-case solution.
The involution-translate observation and defect-vector argument are
reused from the project's earlier F003 audit; the accepted marking
and its diagonal subfield intersection proof are reused unchanged
from \cite{Counterexample}. Proposition~\ref{criterion} assembles
these observations into the corrected matrix criterion.

Wilson's characterization, Suzuki group structure, the quoted
Hubard--Leemans reduction, and the polytope reconstruction theorem
used in the counterexample remain explicit attributed premises.
The source's principal rank theorem and intermediate classification
are false; their disposition is by counterexample, not by recovery
of their conclusions. No classification of every chiral Suzuki
polytope, no formal Lean proof, and no compression factor is claimed.
The binary calculation accompanying (7) checks only that small
matrix identity; general proofs in this note require mathematical
review. This supplement adds no principal-refutation count.

\begin{thebibliography}{5}
\bibitem{Spiga} P. Spiga, \emph{On the rank of Suzuki polytopes:
an answer to Hubard and Leemans}, arXiv:2010.04080v1,
8 October 2020.
\bibitem{HL} I. Hubard and D. Leemans, \emph{Chiral polytopes and
Suzuki simple groups}, Rigidity and Symmetry, Fields Institute
Communications 70 (2014), 155--175.
\bibitem{Suzuki} M. Suzuki, \emph{On a class of doubly transitive
groups}, Annals of Mathematics 75 (1962), 105--145.
\bibitem{Wilson} R. A. Wilson, \emph{A new approach to the Suzuki
groups}, Mathematical Proceedings of the Cambridge Philosophical
Society 148 (2010), 425--428.
\bibitem{Counterexample} Human Compression project,
\emph{A chiral $4$-polytope with automorphism group $\Sz(32768)$},
accepted F003 correction, public proof edition v1,
\texttt{corrections/suzuki-rank/v1/proof}.
\end{thebibliography}
\end{document}
